% SRC: models-lane2/STATUS.md licence tables + blockers; ADR-002; DESIGN.md deployment envelope
\section{Licence and deployment envelope}
\label{sec:licence}

The shipped configuration must fit the UoS target: Windows 11, RTX 5090, \SI{32}{\giga\byte} VRAM,
sm\_120 (Blackwell). Development and validation ran on 48\,GB RTX 6000 Ada cards. Per-stage VRAM
peaks are to be measured under an enforced 32\,GB ceiling before any config is declared 5090-fit,
and a real Windows/5090 acceptance run is the release gate. If that run cannot be executed from the
development container, it is recorded as an explicit release blocker rather than hand-waved.
\todo{Windows/5090 acceptance run on real hardware: this remains the release gate.} The measured
evidence so far is under budget on every path. The object pass's instrumented stage (identify)
peaked at \SI{6.6}{\giga\byte} (Figure~\ref{fig:vram}); generation reserved at most
\SI{9.1}{\giga\byte} on Lane~A and \SI{16.6}{\giga\byte} on Lane~B (Figure~\ref{fig:bakeoff}). All
clear the 32\,GB ceiling, Lane~A with roughly $3.5\times$ headroom, so the acceptance run is
expected to confirm rather than discover the envelope.

\begin{figure}[h]
\centering
% SRC: runs/nc04-object-pass1/summary.json via scripts/make_closeout_charts.py
\includegraphics[width=\columnwidth]{nc04-object-pass1/charts/vram_runtime.pdf}
\caption{nc04 per-stage runtime with the measured VRAM peak against the 32\,GB ship ceiling.
Only \texttt{identify} was VRAM-instrumented in pass~1; the others are marked as unmeasured
rather than assumed.}
\label{fig:vram}
\end{figure}

Licence gating decides what can ship. Table~\ref{tab:licence} reproduces the audited findings. The
decisive point: ReconViaGen v0.5 does not load the Facebook VGGT checkpoint --- it loads
\texttt{Stable-X/vggt-object-v0-1}, a VGGT-architecture derivative with \emph{no declared licence},
so a commercial deployment must treat it as encumbered by Meta's VGGT CC-BY-NC terms until clarified.
The plain VGGT arm loads \texttt{facebook/VGGT-1B} (CC-BY-NC, research/eval only); its commercial
checkpoint is gated behind a manual application form and was not accepted on the owner's behalf.
Consequently \textbf{Lane~B is research/eval-only for this closeout} unless the commercial checkpoint
is obtained. Lane~A (TRELLIS.2-4B, MIT) and everything else in the finishing path are clean. Every
checkpoint used is recorded in lineage: repo, commit, hash and licence class.

\begin{table}[t]
\centering\small
\caption{Model licence gate (audited 2026-08-28).
% SRC: models-lane2/STATUS.md; ADR-002 gate table
}
\label{tab:licence}
\setlength{\tabcolsep}{3pt}
\begin{tabular}{@{}p{0.40\columnwidth}p{0.30\columnwidth}p{0.20\columnwidth}@{}}
\toprule
Component & Licence & Gate \\
\midrule
ReconViaGen v0.5 code & MIT & clean \\
microsoft/TRELLIS.2-4B & MIT & clean \\
Stable-X/trellis-vggt-v0-2 & MIT & clean \\
Stable-X/vggt-object-v0-1 & none declared & treat as NC \\
facebook/VGGT-1B & CC-BY-NC-4.0 & eval only \\
facebook/VGGT-1B-Commercial & gated form & not accepted \\
Hunyuan3D-2.1 & Tencent custom & review before ship \\
SAM3.1 checkpoint (local) & unverified & verify before ship \\
\bottomrule
\end{tabular}
\end{table}

Two environment blockers keep Lane~B unreproduced here. ReconViaGen pins conda py3.10 +
torch 2.4/cu121 with five source-compiled CUDA extensions and cp310 prebuilt wheels
(flash-attn, kaolin, xformers), none installable on the container's nix py3.12; and it ships only a
Gradio demo, so a headless run needs a custom driver written after the environment is solved. The
recommended route is a \texttt{uv}-managed standalone py3.10 in the workspace (mechanism confirmed;
the Docker-via-host-socket route is a documented mount trap). The VGGT arm is fully operational: a
smoke test inferred six gallery frames in \SI{0.94}{\second} at \SI{8.1}{\giga\byte} peak. Lane~A
TRELLIS.2-4B weights (\SI{16}{\giga\byte}) are staged on disk by symlink from a read-only mount, so
no re-download was needed.
